\section{The smallest omitted increments}

For $c=5$ and $c=7$, Theorem~\ref{thm:finite-start} gives the complete candidate ranges
\[
 m<(5+3)(3\cdot5+5)=160,
\]
\[
 m<(7+3)(3\cdot7+5)=260.
\]
Thus a complete exact computation for $1\le m<260$ settles both cases.

The certificate records, for each such $m$, the first absorption index $t$, the resulting increment $c$, and the value $b_t=c(t+1)$. The rows needed to establish the smaller attained increments are shown below.

\begin{table}[ht]
\centering
\caption{Explicit witnesses below the first two omissions.}
\label{tab:witnesses}
\begin{tabular}{@{}rrrr@{}}
\toprule
Eventual increment $c$ & Start $m$ & First absorption index $t$ & $b_t$ \\
\midrule
1 & 3  & 2 & 3 \\
2 & 5  & 7 & 16 \\
3 & 13 & 9 & 30 \\
4 & 19 & 5 & 24 \\
6 & 41 & 7 & 48 \\
\bottomrule
\end{tabular}
\end{table}

\begin{theorem}[Certified nonsurjectivity]
\label{thm:nonsurjective}
The attainable eventual-increment spectrum omits $5$ and $7$, while $1,2,3,4,$ and $6$ occur. Hence $5$ and $7$ are the two smallest omitted positive integers.
\end{theorem}

\begin{proof}
The complete certificate contains every start $1\le m<260$. No row has eventual increment $5$ or $7$. By the finite-start bound, this exhausts every possible witness for $5$ and every possible witness for $7$. Table~\ref{tab:witnesses} gives explicit witnesses for $1,2,3,4,$ and $6$.
\end{proof}
